\chapter{Software Architecture and Engineering}
\label{chap:software_architecture}

\section{System Requirements}
\label{sec:system_requirements}
% TODO: Present functional requirements (encode, generate droplets, convert to DNA, apply screening, simulate losses, decode, collect metrics) and non-functional requirements (modularity, reproducibility, extensibility, readability).

\section{Architectural Decisions}
\label{sec:architectural_decisions}
% TODO: Describe transition from procedural scripting to modular architecture, highlighting separation of concerns, configuration encapsulation, and independent component testing.

\section{Modular Organization}
\label{sec:modular_organization}
% TODO: Detail module structure and responsibilities: config.py, soliton.py, lt_codec.py, screening.py, channel.py, decoder.py, metrics.py, experiments.py, plots.py.

\section{Data Models}
\label{sec:data_models}
% TODO: Formalize entities: segment, droplet, experiment configuration, and experimental result, detailing attributes (seed, degree, payload, DNA sequence, covered blocks, metrics).

\section{Pipeline Flow}
\label{sec:pipeline_flow}
% TODO: Detail sequential flow:
% 1. Message input/definition
% 2. Block segmentation
% 3. LT droplet generation based on RSD
% 4. Seed inclusion for deterministic index reconstruction
% 5. Conversion to DNA
% 6. Biological screening
% 7. Loss/erasure simulation
% 8. Peeling decoding
% 9. Message reconstruction
% 10. Metric consolidation

\section{Implementation Decisions}
\label{sec:implementation_decisions}
% TODO: Document implementation choices: deterministic zero-byte padding, CRC-8 role/handling, PRNG usage separation (index regeneration vs. payload scrambling vs. channel loss events).

\section{Testing and Internal Validation}
\label{sec:testing_and_validation}
% TODO: Document unit tests: bytes-DNA-bytes conversion, seed-based index generation, RSD calculation, screening rules, and lossless recovery validation.
